\section{Automated Hyperparameter Optimization Framework}
\label{sec:hyperparameter_framework}

Identifying the optimal configuration for a GAN is a complex high-dimensional search problem.
To address this, an automated optimization infrastructure was integrated into the pipeline (\texttt{scripts/tune\_hyperparameters.py}), leveraging the \textbf{Optuna} framework~\cite{optuna_2019}.

\subsection{Optimization Strategy}
Instead of relying on inefficient Grid Search or Random Search methods, the framework utilizes a \textbf{Tree-structured Parzen Estimator (TPE)} sampler~\cite{optuna_2019}.
TPE is a Bayesian optimization algorithm that models the probability density function of the objective metric.
As the tuning process progresses, TPE learns which regions of the hyperparameter space are most likely to yield improvements, progressively concentrating the search on the most promising configurations.

\subsection{Search Space and Objective Function}
The optimization is defined over a specific Search Space (e.g., $LR_G \in [10^{-5}, 10^{-3}]$, where $LR_G$ denotes the generator learning rate, $\lambda_{cycle} \in [5, 20]$, where $\lambda_{cycle}$ weights the cycle-consistency loss).
The \textbf{Objective Function} defined for this study drives the TPE sampler to maximize the \textbf{Structural Similarity Index (SSIM)} on the validation set.
SSIM was selected as the target metric because, unlike pixel-wise Mean Squared Error, it correlates strongly with the preservation of structural geological features, guiding the search towards models that maintain stratigraphic integrity rather than just minimizing noise variance.

